% Literature Review & Paper Notes
% Chapter 8 of Amplitudes Notes

\chapter{Literature Review and Paper Notes}
\label{chap:literature-review}

This chapter contains our own notes, breakdowns, and computations related to important papers in the amplitudes literature.

\section{A Brief Introduction to Modern Amplitude Methods}
\label{sec:dixon-2013}

\textbf{Paper:} Lance J. Dixon, ``A brief introduction to modern amplitude methods'', arXiv:1310.5353 (2013).

This is an excellent pedagogical review based on lectures at the 2012 CERN Summer School and TASI 2013. It is one of the best entry points into modern amplitude methods.

\subsection{Key Concepts Covered}

\begin{itemize}
\item \textbf{Color decomposition}: Trace-based decomposition of QCD amplitudes into color-ordered partial amplitudes.
\item \textbf{Spinor helicity formalism}: Review of massless spinors, angle and square brackets, polarization vectors expressed as spinor bilinears.
\item \textbf{Factorization}: Universal soft and collinear factorization of tree-level gauge theory amplitudes.
\item \textbf{Parke-Taylor MHV amplitudes}: The famous one-term formula for maximally helicity violating gluon scattering.
\item \textbf{BCFW recursion}: On-shell recursion relations at tree level.
\item \textbf{Unitarity methods}: Brief introduction to how these ideas extend to loop level.
\end{itemize}

\subsection{Important Results \& Equations}

\textbf{Parke-Taylor MHV formula (color-ordered):}
\begin{equation}
A_n^{\text{tree}}(1^-,2^-;3^+,\dots,n^+) = \frac{\langle 1 2 \rangle^4}{\langle 1 2 \rangle \langle 2 3 \rangle \cdots \langle n 1 \rangle}
\label{eq:parke-taylor}
\end{equation}

\textbf{BCFW recursion relation:}
\begin{equation}
A_n = \sum_{P} A_L(z_P) \frac{1}{P^2(z_P)} A_R(z_P)
\label{eq:bcfw}
\end{equation}
where the sum is over factorization channels and $z_P$ is the value that puts the internal propagator on-shell.

\textbf{Polarization vectors in spinor helicity:}
\begin{align}
\varepsilon^+_\mu(k;q) &= \frac{\langle q | \gamma_\mu | k ]}{\sqrt{2} \langle q k \rangle}, \\
\varepsilon^-_\mu(k;q) &= \frac{[q | \gamma_\mu | k \rangle}{\sqrt{2} [q k ]}.
\end{align}

\subsection{Why This Paper Matters for Our Project}

This paper connects directly to our existing work:
- It builds on the spinor technology from Srednicki Ch. 48 and Ch. 60.
- The Parke-Taylor formula is the seed amplitude we will verify with Cadabra2.
- BCFW recursion is the next major topic after MHV (see Ch. 05 in our notes).

We will add detailed derivations, Cadabra2 checks of the spinor identities, and explicit computations of low-point MHV amplitudes in subsequent sections.

\newpage

\section{Single-Minus Gluon Tree Amplitudes Are Nonzero}
\label{sec:smga-2025}

\textbf{Paper:} Guevara, Lupsasca, Skinner, Strominger, Weil (on behalf of OpenAI), ``Single-minus gluon tree amplitudes are nonzero'', arXiv:2602.12176v2 (2025).

\textbf{Significance:} This is the first physics discovery made by an AI system (GPT-5.2 Pro). The key formula was \emph{conjectured} by GPT and then \emph{proved} by a new internal OpenAI model.

\subsection{The Discovery}

It was long believed that single-minus tree-level gluon amplitudes $\mathcal{A}_n(1^-,2^+,\dots,n^+)$ vanish at tree level for generic kinematics. The standard argument uses a clever choice of reference spinor to show all polarization vectors are orthogonal, so the amplitude requires more powers of momenta in the numerator than vertices can supply.

The \textbf{loophole}: when all angle brackets vanish ($\langle ij\rangle = 0$ for all $i,j$), the reference spinor choice becomes singular. This ``half-collinear'' regime exists in $(2,2)$ Klein signature or for complexified momenta.

\subsection{Key Equations}

\textbf{The stripped amplitude ansatz:}
\begin{equation}
\mathcal{A}_n = i^{2-n} A_{1\cdots n} \prod_{a=2}^{n} \delta(\langle 1a\rangle) \cdot \delta^2\!\left(\sum_{i=1}^{n} \langle ri\rangle \tilde\lambda_i\right)
\end{equation}

\textbf{The main result (Eq.~16 of the paper):}
\begin{equation}
A_{1\cdots n}\Big|_{\mathcal{R}_1} = \frac{1}{2^{n-2}} \prod_{m=2}^{n-1} \left(\mathrm{sg}_{m,m+1} + \mathrm{sg}_{1,2\cdots m}\right)
\label{eq:smga-main}
\end{equation}
where $\mathrm{sg}_{ij} = \mathrm{sg}([\tilde\lambda_i \tilde\lambda_j])$ and $\mathcal{R}_1$ is the kinematic region where particle 1 has $\omega_1 < 0$ and all others have $\omega_a > 0$.

\textbf{Concrete low-point examples:}
\begin{align}
A_{123}\big|_{\mathcal{R}_1} &= \tfrac{1}{2}(\mathrm{sg}_{23} + \mathrm{sg}_{12}), \\
A_{1234}\big|_{\mathcal{R}_1} &= \tfrac{1}{4}(\mathrm{sg}_{23} + \mathrm{sg}_{12})(\mathrm{sg}_{34} + \mathrm{sg}_{41}).
\end{align}

Each factor is $\frac{1}{2}(\pm 1 \pm 1) \in \{-1, 0, +1\}$, so the stripped amplitude is \textbf{piecewise-constant}.

\subsection{Proof Strategy}

The proof proceeds in three steps:
\begin{enumerate}
\item Show that in $\mathcal{R}_1$, the vertex function $V_{\tilde\lambda_2\cdots\tilde\lambda_n} = 0$ (a causality condition).
\item Show the Berends--Giele recursion collapses: $A_{1\cdots n} = \bar{V}_{\tilde\lambda_2\cdots\tilde\lambda_n}$.
\item Evaluate $\bar{V}$ using momentum conservation to recover Eq.~\eqref{eq:smga-main}.
\end{enumerate}

A key ingredient is the ``master identity'' (App.~A), a generalization of a standard distributional identity involving sign functions and $\delta$-functions.

\subsection{Open Directions (from the paper)}

The authors note several extensions we should explore:
\begin{itemize}
\item General formula outside $\mathcal{R}_1$ (they give explicit results through $n=6$)
\item Graviton amplitude generalization
\item Supersymmetric extensions
\item Connection to $\mathcal{L}w_{1+\infty}$ algebra and celestial holography
\item Mellin transforms yielding Lauricella functions
\item Possible yet-simpler expressions via different variable choices
\end{itemize}

\subsection{Why This Paper Matters for Our Project}

This paper demonstrates that:
\begin{enumerate}
\item AI systems can make genuine discoveries in scattering amplitudes.
\item The discovery method was pattern recognition from computed data (GPT conjectured the formula from $n=3,\ldots,6$ examples).
\item The proof was also AI-assisted.
\end{enumerate}

Our goal is to reproduce this methodology and find new results in related areas.

\newpage
